\section{Introduction}

\subsection{Topic and Problem Statement}

The recently passed Regulation 2025/2818 laying down additional procedural rules on the enforcement of Regulation (EU) 2016/679 (hereinafter ‘Procedural Regulation’)\autocite{ProceduralRegulation} seeks to further proceduralize cross-border enforcement of the General Data Protection Regulation (GDPR)\autocite{GDPR}. 

The GDPR\autocite{GDPR} is the cornerstone of the European Union's (EU) so-called Digital Acquis.\autocite[223]{zanfirGDPREDPS} Entered into force in 2018, and replacing the Data Protection Directive of 1995, it has often been heralded as the most comprehensive package of data protection and privacy regulations.\autocite[498]{mustertEuropeanUnionProceduralising2026} 

The GDPR aims to provide a consistent and harmonized framework to protect the fundamental right to data protection across the EU\autocite[recital 7]{GDPR}, as enshrined in the Charter of Fundamental Rights (CFR)\autocite[8(1)]{CFR} and in the Treaty on the Functioning of the European Union (TFEU)\autocite[16(1)]{TFEU}. As recognized by the Regulation itself, the need for predictable and effective protection of personal data has become incresingly acute, given the scale at which information is being processed.\autocite[recital 6]{GDPR}

While the choice of legislative instrument and the language of the GDPR highlight the need for a strong framework for data protection, which has become inseparable from the exercise of other fundamental rights in democratic societies, the enforcement of the GDPR has encountered practical challenges, especially in the context of cross-border cases. The Procedural Regulation seeks to address these challenges by further defining the steps involved in cooperation between Member States' (MS) authorities, introducing earlier consensus-finding and time limits for every new stage of the procedure.

This thesis seeks to analyze the Procedural Regulation, and its potential impact on the enforcement of the GDPR in cross-border cases. It will do so in light of the general principle of good administration, with a view to investigate whether the new procedural rules will lead to a more effective and consistent enforcement of the GDPR, contributing to the protection of fundamental rights across the EU.